\owner{Member 1 integrates; Members 2--4 supply pair-specific literature notes.}
\subsection{Motivation}
Training a neural network requires minimizing an empirical loss over its parameters.
Stochastic gradients make this feasible, but the trajectory depends on step sizes, gradient
scaling, and the information retained from earlier updates \citep{bottou2018}.
A comparison based only on final classification accuracy can miss slow optimization,
unstable updates, or a large tuning cost. Our objective is therefore to connect differences
in training behavior to specific algorithmic mechanisms.

A familiar dataset is useful when the experiment isolates a meaningful factor.
For example, \citet{zhang2024hessian} investigate block heterogeneity using quadratic problems
and neural networks, while \citet{xie2025geometry} examine the effect of coordinate geometry.
These studies motivate controlled diagnostics alongside real tasks; they do not imply that
classification difficulty alone measures poor conditioning.
\subsection{Related work}
We organize the literature around four modifications: dynamic distance-based step sizes
\citep{ivgi2023dog}, the relationship between scheduling and iterate averaging
\citep{defazio2024schedule}, changes to adaptive normalization \citep{taniguchi2024adopt},
and the use of two gradient moving averages \citep{pagliardini2025ademamix}.
Each modification addresses a different question, so each is paired with an appropriate
classical baseline. Our study uses small networks and CPU budgets; it does not reproduce
the scale of the original papers or claim that all four modifications solve the same problem.
\subsection{Research questions and scope}
We ask: (RQ1) Does each modification improve loss reduction or reduce tuning dependence
under a shared budget? (RQ2) Does a targeted diagnostic support the proposed mechanism?
(RQ3) Do observations persist across image classification and character prediction, and at
what cost? Each member will independently implement and evaluate one pair. The planned
contribution is a reproducible empirical comparison and an explanation of observed trade-offs,
not a new optimizer. \fillin{Member 1: update scope to match the completed experiments.}
